\section{Introduction}\label{sec:intro}

Retrieval-augmented generation (RAG) answers a question in two steps. A retriever selects passages from a corpus, and a language model, the reader, writes an answer from them \citep{lewis2020rag,izacard2021fid,ram2023incontext}. Teams that build such systems must choose the retriever. For dense retrieval this means choosing a text embedding model, and the usual guide is a retrieval benchmark such as BEIR \citep{thakur2021beir} or MTEB \citep{muennighoff2023mteb}. These benchmarks score ranked lists against relevance labels. They do not show how much of a retrieval gain survives the reader.

This question matters most for teams with limited compute. A larger encoder costs more to run over a corpus and adds latency to every query. If a retrieval gain turns into a much smaller answer gain, a cheaper encoder may serve as well. It is already known that better retrieval does not reliably give better answers \citep{salemi2024erag,mathur2026policy,medrano2026fusion,ghosh2026evidence}, that retrieval metrics can still track generation quality when the two are aligned \citep{samuel2026beyond,hitl2026comparative}, and that passages which look relevant but lack the answer can lower accuracy \citep{cuconasu2024power,amiraz2025distracting,shi2023distracted}. What is less clear is how large the answer gain is, relative to the retrieval gain, when one interchangeable encoder replaces another, and where in the query set that gain appears.

Earlier studies relate retrieval quality to answer quality through correlations across systems or topics \citep{salemi2024erag,samuel2026beyond,trappolini2025udcg}, through a ratio of answer F1 to recall across RAG methods \citep{peng2026neocorrag}, or through manipulations of a single retriever \citep{leto2024toward}. Correlations describe whether systems are ranked in the same order; they do not say how many answer points one retrieval point buys. Two encoders can be ranked correctly by Hit@5 while the answer gain between them is a quarter of the retrieval gain. We found no study that reports, for a set of interchangeable open encoders, the paired ratio of answer gain to retrieval gain with its uncertainty, per reader and dataset, next to the measured retrieval-stage cost of each encoder. This study provides such estimates for one controlled setting.

We define a transfer ratio $T(R)$: the gain in exact match (EM) of retriever $R$ over a small baseline encoder, divided by its gain in Hit@5, both measured on the same queries. A ratio of 1 means that EM rises by as many points as Hit@5, and a ratio near 0 means that the retrieval gain does not reach the answer. $T(R)$ is a descriptive ratio of paired mean gains. It is not a probability, not a causal conversion efficiency, and not the share of newly retrieved passages that the reader uses; it depends on the baseline and on the retrieval metric, so it describes a pair of retrievers rather than one encoder; and 1 is a reference point, not a target. We therefore also account for the gain query by query. We estimate $T(R)$ for BM25 and five dense encoders against MiniLM-L6, with two quantised 7--8B readers, on a single-hop dataset (Natural Questions) and a multi-hop dataset (HotpotQA). All runs used free Kaggle T4 GPUs.

We ask four research questions.
\begin{description}
\item[RQ1] How much of a Hit@5 gain over a small encoder reaches EM, and does this depend on the reader and the dataset?
\item[RQ2] Do larger or newer encoders give higher EM than a mid-size encoder (e5-base-v2)?
\item[RQ3] How do context choices change the answer: the number of passages, the kind of non-gold passages, the position of the gold passage, and reranking?
\item[RQ4] What is each encoder's retrieval-stage cost on free hardware, and do our sub-corpus and exact search distort the comparison?
\end{description}

The study ran in three phases. Phase~1 covered one reader on NQ. Before Phase~2, we registered six predictions (P1--P6) in a dated design document and then ran the new cells and experiments. Before Phase~3, we registered two more (P7, P8) and ran both readers on retrieval from the full NQ corpus, to test whether the 300,000-passage sub-corpus that all other answer results use had changed the NQ ratios. We report whether each stated criterion was met, including the three that were not met in at least one condition.

The study makes four contributions.
\begin{itemize}
\item Paired estimates of the answer gain per unit of retrieval gain for six retrievers, two readers and two datasets, with bootstrap intervals and a reporting rule fixed in advance, their sensitivity to the answer metric, the baseline, the retrieval metric and the number of passages read, and a registered replication of the NQ ratios on the full corpus (Sections~\ref{sec:res-transfer}--\ref{sec:res-full}).
\item Evidence that, for multi-hop questions, the size of the transfer depends on whether retrieval is scored by complete evidence: with Both@5 instead of Hit@5 the HotpotQA ratios fall to the NQ range (Section~\ref{sec:res-sens}).
\item A query-level accounting that splits each EM gain exactly into queries that gained, kept or lost labelled evidence, which shows where the NQ and HotpotQA gains sit, and an answer-string check of the passages behind NQ retrieval misses (Section~\ref{sec:res-decomp}).
\item Measured retrieval-stage cost of each encoder on a free GPU, controlled checks of context composition, reranking, approximate search and the sub-corpus, and a full record of every deviation and failed run (Sections~\ref{sec:res-context}--\ref{sec:res-cost} and~\ref{sec:deviations}).
\end{itemize}

Section~\ref{sec:related} reviews related work. Sections~\ref{sec:method} and~\ref{sec:setup} describe the method and setup. Section~\ref{sec:results} reports the results, Section~\ref{sec:discussion} discusses them, and Section~\ref{sec:limits} lists the limits of the study.
